\documentclass[../HDF5_RFC.tex]{subfiles}

\begin{document}

\section{Introduction}
\label{sec:intro}

\subsection{Purpose}

This document describes the high-level design for a new string-based interface
for configuring HDF5 I/O filters. The feature adds a usability layer on top of
the existing filter machinery, allowing users to specify filter names and
parameters as human-readable strings rather than constructing raw \texttt{cd\_values}
arrays of \texttt{unsigned int}.

\subsection{Scope}

This design covers:

\begin{itemize}
\item New public API functions in the \texttt{H5P} and \texttt{H5Z} modules for string-based
  filter configuration.
\item A new version of the filter plugin class (\texttt{H5Z\_class3\_t}) with optional
  string-configuration callbacks.
\item A canonical name registry for filter lookup by name.
\item Extensions to the plugin loading infrastructure (\texttt{H5PL}) for name-based
  plugin discovery.
\item Integration points for HDF5 command-line tools (\texttt{h5repack}, \texttt{h5dump}).
\item The Fortran and Java bindings that ship with HDF5 (3 \texttt{H5P} subroutines and
  4 \texttt{H5Z} subroutines for Fortran; corresponding Java wrappers).
\end{itemize}

This design does not cover:

\begin{itemize}
\item Changes to the on-disk format or object header layout.
\item Changes to the filter execution pipeline itself.
\item External high-level language binding implementations (h5py, etc.), though
  the design is motivated in part by their needs.
\item Data-attribute-dependent compression strategies (e.g., choosing different
  filters based on dimensionality or datatype), which could be layered on top
  of this feature in the future.
\end{itemize}


\section{Motivation}
\label{sec:motivation}

\subsection{Problem Statement}

The current generic interface for configuring HDF5 filters is \texttt{H5Pset\_filter},
which accepts an array of \texttt{unsigned int} values (\texttt{cd\_values}). This interface
presents three categories of problems.

\verysubsection{Type punning}
Filters that accept floating-point parameters (rates,
tolerances, precision thresholds) require users to manually bit-cast \texttt{double}
or \texttt{float} values into one or more \texttt{unsigned int} slots. This is error-prone,
endian-sensitive, and difficult or impossible to express naturally in
high-level languages such as Python, Java, and MATLAB without dedicated
wrapper code. For example, configuring H5Z-ZFP in rate mode currently
requires:

\begin{minted}{c}
unsigned int cd_vals[4];
double rate = 3.5;
cd_vals[0] = H5Z_ZFP_MODE_RATE;
memcpy(&cd_vals[2], &rate, sizeof(double));
H5Pset_filter(plist, H5Z_FILTER_ZFP, flags, 4, cd_vals);
\end{minted}

\verysubsection{CLI usability}
Command-line tools such as \texttt{h5repack} require users to pass
opaque integer sequences (e.g., \texttt{-f UD=32013,0,4,1,0,0,1074921472}). External
helper scripts exist solely to translate human-readable settings into these
sequences.

\verysubsection{Opaque introspection}
When examining a file's dataset creation properties,
the \texttt{cd\_values} array is meaningless without external documentation. There is
no standard way to recover the original intent of the configuration (e.g.,
``ZFP rate mode, rate=3.5'').

\subsection{Precedent Within HDF5}

HDF5 1.14 introduced string-based configuration for Virtual File Drivers:

\begin{minted}{c}
H5Pset_driver_by_name(fapl, "sec2", NULL)
\end{minted}

This API allows drivers to be selected and configured by name at runtime
without requiring compiled-in constants. The VOL subsystem provides analogous
by-name functionality through \texttt{H5VLregister\_connector\_by\_name()} and
\texttt{H5VLget\_connector\_id\_by\_name()}, though it does not have a direct
\texttt{H5Pset\_vol\_by\_name} equivalent. The proposed filter configuration API follows
the VFD pattern most closely and extends it with structured parameter passing.

\subsection{Goals}

\begin{enumerate}
\item Users MUST be able to configure any filter using human-readable named
   parameters.
\item Filter plugins MUST be addressable by a canonical name, without requiring
   compiled-in integer constants. (See the open design decision in
   \SectionRef{sec:set-filter-by-name} regarding whether name-based
   identification justifies its added complexity over an ID-based approach.)
\item The on-disk format MUST NOT change.
\item All existing APIs, plugins, and call sites MUST continue to work without
   modification.
\item The design MUST support both built-in and dynamically loaded plugin
   filters.
\item The API MUST be consistent with the structure and conventions of existing
   HDF5 public APIs.
\end{enumerate}

\end{document}
